% chapter_02_PS.tex -- JitJet chapter 02: the parton shower. \input by doc/JitJet.tex (Rule D).
% Every number is from test/02_PS/out/ (ledger_02_PS.txt, genPartAnalyzer_02_PS.log, step1_GEN_02_PS.log) or
% test/01_ME/out/ and test/04_MPI/out/ for the comparisons; the companion doc/chapter/chapter_02_PS.md lists them in full.
\chapter{02 -- PS: the parton shower dresses the $b$ quark with gluons}
\label{ch:02_PS}

\section{What this chapter is about}

In chapter 01 the event was the bare matrix element: an antiquark and a quark in, a $b$ and a $\bar b$ out, each with
$\pt = 113.2\GeV$, fourteen entries in the record and nothing else. A quark that is kicked to $113\GeV$ does not stay
alone: it radiates gluons, as an accelerated electron radiates photons, and the gluons radiate further gluons or split
into quark pairs. Pythia8 computes this radiation with its \emph{parton shower}: a chain of $1\to2$ branchings
($q\to qg$, $g\to gg$, $g\to q\bar q$) ordered in the transverse momentum of the emission, started at the scale of the
hard process and stopped at a cut-off of about half a GeV, below which perturbation theory is no longer trusted and
the hadronisation of chapter 05 takes over. There are two showers. The \emph{initial-state radiation} (ISR) is
emitted by the incoming $\bar u$ and $u$ before they collide and is generated backwards from the hard process towards
the protons, guided by the parton distribution functions. The \emph{final-state radiation} (FSR) is emitted by the
outgoing $b$ and $\bar b$ after the collision. This chapter switches on exactly these two steps,
\code{PartonLevel:ISR} and \code{PartonLevel:FSR}, and leaves multiparton interactions, hadronisation and decays off,
so that the final state is still made of partons.

What happens to our event is small in size and large in meaning. The record grows from 14 to 34 entries; the shower
branches five times, all on the $b$ side: one $b\to bg$, three $g\to gg$ and one $g\to d\bar d$. The $b$ quark that
leaves the shower carries $46.9\GeV$ instead of $113.2\GeV$; the six partons of its system together still carry
$113.6\GeV$, all within $\Delta R = 0.31$ of the quark. The $\bar b$ does not radiate at all and the ISR evolution
produces no emission. The first lesson of the journey is written in these numbers: the momentum of the hard process
is carried by the jet, not by the quark, and the jet exists before any hadron does.

\begin{plainwords}[The parton shower in plain words]
The hard process is calculated at one scale, $\ptHat = 113\GeV$. Between that scale and the hadron scale of
$1\GeV$ there are two orders of magnitude in which the strong coupling is still perturbative, and where the
probability of emitting a gluon of transverse momentum $\pt$ grows like $\alpha_s\,\mathrm{d}\pt^2/\pt^2$: soft and
collinear gluons are emitted copiously. The shower generates them one by one, from the hardest down, each emission
taking its momentum from the parton that radiates and from a \emph{recoiler}, the colour partner of the radiator
(a dipole). Every emission is kept in the record, together with a new copy of every parton whose momentum changed.
Nothing is created or lost: the energy of the event stays $13600\GeV$, the momentum of the $b$ is shared among the
$b$ and its radiation.
\end{plainwords}

\section{From the previous chapter}

The input of this chapter is the configuration of chapter 01 with two switches turned on: the only differences
between \code{test/01\_ME/step1\_GEN\_01\_ME\_cfg.py} and \code{test/02\_PS/step1\_GEN\_02\_PS\_cfg.py} are
\code{PartonLevel:ISR = on}, \code{PartonLevel:FSR = on} and the file names. The same fragment, the same tune, the
same random seed (\code{initialSeed = 123456789} of the \code{generator} engine in
\code{IOMC/RandomEngine/python/IOMC\_cff.py}) and the same beam energy mean that Pythia8 draws the same hard process:
the incoming $\bar u$ with $x_1 = 2.174\times10^{-3}$ and $u$ with $x_2 = 0.3901$, $\sqrt{\hat s} = 396.1\GeV$,
$\ptHat = 113.2\GeV$, the $b$ at $\eta = -1.44$, $\phi = -2.15$ and the $\bar b$ at $\eta = -3.75$, $\phi = 0.99$ are
identical in the two logs and ledgers, and the first page of the analyzer PDF, the display of the hard process, is
the same picture as in chapter 01.

What changes is everything that comes after the hard process. The shower adds 15 entries with the status codes
51--53 and raises the final state from 6 to 11 partons. The beam remnants, which Pythia8 builds only after the
showers, are drawn with a different random sequence: in chapter 01 the proton of the $\bar u$ left a $ud_1$ diquark,
a $u$ of $132.8\GeV$ and a $u$ of $9.1\GeV$ at $\eta = 3.36$; here it leaves a $ud_0$ diquark and two $u$ quarks of
$415\GeV$ and $552\GeV$ beyond $\abseta = 7$. This is the first sign of a caveat that matters for the whole
book, stated plainly in section~\ref{sec:02_PS_numbers}: the hard process is the same event in chapters 01 to 06,
the emissions of the shower are not.

\section{The numbers of the event}
\label{sec:02_PS_numbers}

\paragraph{The shower that Pythia8 ran.} Table~\ref{tab:02_PS_settings} lists the settings that govern the two
showers, from the settings table that Pythia 8.311 prints at the start of \code{step1\_GEN\_02\_PS.log} (it prints
only the parameters that differ from its defaults) and from the Pythia8 documentation for the rest. The shower
model is the default \emph{simple shower} of Pythia8 (\code{PartonShowers:model = 1}; the FSR class
\code{SimpleTimeShower} is named in the warning of the log, its ISR counterpart is \code{SimpleSpaceShower}), a \pt-ordered dipole shower:
emissions are generated in decreasing \pt, each from a radiator--recoiler pair that shares a colour line. The CP5
tune sets the strong coupling of both showers to $\alpha_s(M_Z) = 0.118$ with two-loop running, the value of the
PDF set \code{NNPDF31\_nnlo\_as\_0118}, instead of the Pythia8 default $0.1365$ with one-loop running; a smaller
coupling means fewer and softer emissions. The starting scale of both showers is the factorisation scale of the
hard process, $\sqrt{Q^2} = 113.3\GeV$ (the Info listing of the log: \code{Max pT scale for ISR = 1.133e+02, FSR =
1.133e+02}). The FSR stops at \code{TimeShower:pTmin}; its default $0.5\GeV$ is too low for the two-loop coupling and
Pythia8 raised it itself: \code{Warning in SimpleTimeShower::init: pTmin too low, raised to 0.611}. The ISR is
regularised by \code{SpaceShower:pT0Ref = 2.0} ($2.0\GeV$), a Pythia8 default that the tune does not touch.

\begin{table}[htbp]
\centering\footnotesize
\caption{The settings of the two showers in this run, from the Pythia8 settings table of
\code{test/02\_PS/out/step1\_GEN\_02\_PS.log} (``Now'' and ``Default'' columns) and, for the parameters left at
their default, from the Pythia 8.311 documentation (\code{share/Pythia8/xmldoc}). ``CP5'' is
\code{Configuration/Generator/python/MCTunes2017/PythiaCP5Settings\_cfi.py}.}
\label{tab:02_PS_settings}
\begin{tabular}{lll}
\toprule
parameter & this run (set by) & Pythia8 default \\
\midrule
\code{PartonLevel:ISR}, \code{PartonLevel:FSR} & on, on (\code{run\_02\_PS.sh}) & on, on \\
\code{PartonLevel:MPI}, \code{HadronLevel:all} & off, off (\code{run\_02\_PS.sh}) & on, on \\
\code{PartonShowers:model} & 1, the simple shower (default) & 1 \\
\code{TimeShower:alphaSvalue}, \code{alphaSorder} & 0.118, 2 (CP5) & 0.1365, 1 \\
\code{SpaceShower:alphaSvalue}, \code{alphaSorder} & 0.118, 2 (CP5) & 0.1365, 1 \\
\code{TimeShower:pTmin} & $0.5\GeV$, raised to $0.611\GeV$ at init & $0.5\GeV$ \\
\code{SpaceShower:pT0Ref}, \code{ecmPow} & $2.0\GeV$, 0 (default) & $2.0\GeV$, 0 \\
\code{PDF:pSet} & \code{LHAPDF6:NNPDF31\_nnlo\_as\_0118} (CP5) & 13, NNPDF2.3 LO \\
starting scale of ISR and FSR (Info listing) & $113.3\GeV = \sqrt{Q^2}$ of the hard process & \\
\bottomrule
\end{tabular}
\end{table}

\paragraph{How to read the record.} Pythia8 marks every entry of the event record with a status code that says in
which step it was produced; CMSSW copies the code into the \code{genParticles} collection. The hard process uses
21 (incoming) and 23 (outgoing); the ISR uses 41--49 (41 incoming on the space-like branch, 43 outgoing produced by a
branching, 44 outgoing shifted by a branching); the FSR uses 51--59: 51 is an outgoing parton produced by a branching,
52 is a \emph{copy} of a parton whose momentum changed because it was the recoiler of a branching, 53 the same for
an incoming parton. A parton is therefore never modified in place: when the $b$ radiates, the record keeps the
original $b$ (status 23), adds the $b$ and the gluon after the branching (status 51) and, each time the $b$ is the
recoiler of a later branching, a new copy (status 52). The analyzer follows the ``last copy'' of the $b$ through the
chain. The final-state partons are copied once more with status 1 after the beam remnants are attached (the
primordial transverse momentum of the incoming partons is added at that step: the $\bar b$ goes from $113.221$ to
$114.048\GeV$ without having radiated).

\begin{table}[htbp]
\centering\footnotesize
\caption{The five FSR branchings of the event, read from the \code{genParticles} of
\code{test/02\_PS/out/step1\_GEN\_02\_PS.root} with FWLite (the full record is in the companion). Index = position
in the collection; the recoiler is the colour partner that absorbs the recoil and is copied with status 52 (outgoing)
or 53 (incoming). The last column is the vector sum of the \pt{} of the two daughters.}
\label{tab:02_PS_chain}
\begin{tabular}{@{}c@{\hspace{6pt}}l@{\hspace{6pt}}r@{\hspace{6pt}}l@{\hspace{6pt}}r@{\hspace{6pt}}l@{\hspace{6pt}}r@{}}
\toprule
branching & radiator (idx, st) & $\pt$ [GeV] & daughters (idx, st) & $\pt$ [GeV] & recoiler (idx $\to$ copy, st) & $|\sum\vec\pt|$ \\
\midrule
1 & $b$ (5, 23) & 113.22 & $b$ (6, 51), $g$ (7, 51) & 47.70, 65.58 & incoming $u$ (3 $\to$ 8, 53) & 113.22 \\
2 & $g$ (7, 51) & 65.58 & $g$ (9, 51), $g$ (10, 51) & 61.70, 4.11 & incoming $u$ (8 $\to$ 11, 53) & 65.68 \\
3 & $g$ (10, 51) & 4.11 & $d$ (12, 51), $\bar d$ (13, 51) & 2.46, 7.66 & $g$ (9 $\to$ 14, 52) & 10.10 \\
4 & $g$ (14, 52) & 55.60 & $g$ (15, 51), $g$ (16, 51) & 51.94, 5.02 & $\bar d$ (13 $\to$ 17, 52) & 56.95 \\
5 & $g$ (15, 51) & 51.94 & $g$ (18, 51), $g$ (19, 51) & 43.12, 9.66 & $b$ (6 $\to$ 20, 52) & 52.76 \\
\bottomrule
\end{tabular}
\end{table}

\paragraph{The five branchings.} Table~\ref{tab:02_PS_chain} is the shower of this event. The Info listing of the
log counts it: \code{Number of MPI = 1, ISR = 0, FSRproc = 5} (the one ``MPI'' is the hard process itself). The
first branching is the hardest, as the \pt{} ordering demands: the $b$ of $113.22\GeV$ emits a gluon of $65.58\GeV$
and keeps $47.70\GeV$; the two together have the \pt{} of the original $b$ to the last digit and an invariant mass
of $11.6\GeV$, and they are $\Delta R = 0.156$ apart. This is a collinear emission, not a wide-angle one: the gluon
carries more \pt{} than the quark but flies in the same direction. The gluon then splits into two gluons, the softer
of which splits into a $d\bar d$ pair, and the harder gluon splits twice more. The recoil of the first two branchings
is taken by the incoming $u$, the colour partner of the $b$ in the $u\bar u\to b\bar b$ colour flow: its energy
moves from $2652.93$ to $2656.06\GeV$, that is $x_2$ from $0.39014$ to $0.39060$ (the two status-53 entries). The
other three recoils go to the shower partons themselves (status 52). The $b$ is recoiler of the last branching and
ends at $46.88\GeV$ (index 20, status 52, the copy the analyzer reports) and at $47.04\GeV$ in its status-1 copy.

\paragraph{Where the momentum went.} The six final-state partons of the $b$ side ($b$, three gluons, $d$, $\bar d$:
indices 24 to 29) have a vector \pt{} of $113.56\GeV$, a scalar sum of $113.75\GeV$, an energy of $256.9\GeV$ and
an invariant mass of $13.8\GeV$; the farthest of them, the $\bar d$ of $6.3\GeV$, is $\Delta R = 0.311$ from the
$b$, the gluon of $43.2\GeV$ is at $0.161$. A jet of radius 0.4 drawn around the $b$ would contain them all and
measure $113.6\GeV$: the shower did not change the jet, it changed what the jet is made of. The analyzer sees it
from the other side: its ``net $b+\bar b$'' is the vector sum of the two \emph{quarks} and jumps from
$0.0$ to $66.4\GeV$, with a scalar sum of $160.1\GeV$ instead of $226.4\GeV$. That imbalance is the radiation
that is still in the record, six entries lower; the vector sum of the $b$ system and the $\bar b$ is $0.6\GeV$, the
rest being the primordial $k_\mathrm{T}$ and the recoil of the incoming $u$. The $\bar b$ at $\eta = -3.75$ radiated
nothing: in a shower every emission is a random draw against the Sudakov factor, and a $b$ quark of $113\GeV$ has a
sizeable probability to emit nothing above the cut-off. The same holds for the ISR: the evolution from $113.3\GeV$
down to the cut-off happened to produce no emission. One reason it is not rare here is kinematic: the $u$ carries
$x_2 = 0.39$, and the backward evolution of a valence quark that already has a large $x$ must find a parent with
an even larger $x$, which the steeply falling parton distribution disfavours.

\begin{plainwords}[Why this shower is not the shower of chapters 04 to 06]
Pythia8 does not run ISR, FSR and multiparton interactions one after the other; it runs them \emph{interleaved} in
one \pt-ordered evolution, in which at every step the next emission is the one with the highest \pt{} among all
three. With MPI switched off, as here, the sequence of random numbers after the hard process is different from the
sequence with MPI on, and so are the emissions. The same event in chapter 04 (\code{test/04\_MPI/out/},
MPI on) has \code{ISR = 28, FSRproc = 64} branchings counted in its log, and its analyzer reports at stage 2 a $b$ of
$91.6\GeV$ and a $\bar b$ of $118.6\GeV$, both status 52, net $27.3\GeV$, against $46.9\GeV$, $113.2\GeV$ (status 23)
and $66.4\GeV$ here. The hard process is the same event, the shower emissions are not identical between this
chapter and the later ones; chapters 04, 05 and 06 are identical to each other at parton level. This chapter shows
what a shower does, isolated; the shower that the detector will see is the one of chapter 04.
\end{plainwords}

\section{Ledger and display}

\begin{table}[htbp]
\centering\small
\caption{The ledger of the stage: the generator record before (\code{test/01\_ME/out/ledger\_01\_ME.txt},
\code{genPartAnalyzer\_01\_ME.log}) and after the shower (\code{test/02\_PS/out/ledger\_02\_PS.txt},
\code{genPartAnalyzer\_02\_PS.log}). $E$ and \pt{} are sums over the final-state (status 1) particles of the region;
the whole event is $13600\GeV$ in both. The $b$ and $\bar b$ rows are the last quark copies the analyzer follows
(status 23 before, status 52 and 23 after); the $b$ system row is the $b$ with its five radiated partons.}
\label{tab:02_PS_ledger}
\begin{tabular}{lrrr}
\toprule
quantity & before (01, ME) & after (02, PS) & difference \\
\midrule
entries in \code{genParticles} & 14 & 34 & $+20$ \\
\quad hard process (status 21--29) & 4 & 4 & 0 \\
\quad FSR (status 51--59) & 0 & 15 & $+15$ \\
\quad ISR (status 41--49) & 0 & 0 & 0 \\
\quad beam remnants (status 61--69) & 2 & 2 & 0 \\
\quad final state (status 1) & 6 & 11 & $+5$ \\
\midrule
final-state energy, all [GeV] & 13600.0 & 13600.0 & 0.0 \\
final-state scalar \pt, all [GeV] & 228.6 & 230.7 & $+2.1$ \\
$\abseta < 2.5$: $n$ / $E$ / \pt{} [GeV] & 1 / 253.1 / 113.6 & 6 / 256.9 / 113.7 & $+5$ / $+3.8$ / $+0.1$ \\
$3.0 < \abseta < 5.2$: $n$ / $E$ / \pt{} [GeV] & 2 / 2423.7 / 113.0 & 1 / 2413.9 / 114.0 & $-1$ / $-9.8$ / $+1.0$ \\
$\abseta > 5.5$: $n$ / $E$ / \pt{} [GeV] & 3 / 10923.1 / 2.1 & 4 / 10929.2 / 2.9 & $+1$ / $+6.1$ / $+0.8$ \\
visible ($\abseta < 5.2$): $n$ / $E$ [GeV] & 3 / 2676.9 & 7 / 2670.8 & $+4$ / $-6.1$ \\
\midrule
$b$ quark \pt{} [GeV] & 113.2 & 46.9 & $-66.3$ \\
$\bar b$ quark \pt{} [GeV] & 113.2 & 113.2 & 0.0 \\
$b$ system ($b$ + radiation) vector \pt{} [GeV] & 113.2 & 113.6 & $+0.4$ \\
net $|\vec\pt(b) + \vec\pt(\bar b)|$ [GeV] & 0.0 & 66.4 & $+66.4$ \\
scalar $\pt(b) + \pt(\bar b)$ [GeV] & 226.4 & 160.1 & $-66.3$ \\
\bottomrule
\end{tabular}
\end{table}

Table~\ref{tab:02_PS_ledger} closes the ledger and figure~\ref{fig:02_PS_display} shows the result. The shower added twenty entries and five final-state partons, all of
them in the tracker region $\abseta < 2.5$ next to the $b$; it moved no energy between the regions beyond what the
re-drawn beam remnants did (the $u$ of $9.1\GeV$ at $\eta = 3.36$ of chapter 01 is gone, the forward region now
holds the $\bar b$ alone; the remnants beyond $\abseta = 5.5$ gained one entry and $6\GeV$). The energy of the
event is conserved to the first decimal and the visible energy changes by $6\GeV$ out of $2671$. The one large number
of the table, $-66.3\GeV$ on the $b$ quark, is not a loss: the $b$ system row shows it was moved to the five partons
that now share the jet with the quark.

\begin{figure}[!tb]
\centering
\includegraphics[width=\textwidth]{02_PS/out/genPartAnalyzer_02_PS_page-2.png}
\caption{The $\eta$--$\phi$ display of the event after the shower (\code{GenPartAnalyzer}, page 2 of
\code{test/02\_PS/out/genPartAnalyzer\_02\_PS.pdf} as \code{genPartAnalyzer\_02\_PS\_page-2.png}; the marker size
grows with $\log\pt$; page 1, the hard process, is the display of chapter 01): the $b$ at $46.9\GeV$ with its five
radiated partons (green triangles, labelled when above $5\GeV$), the untouched $\bar b$, and the vector sum
$b+\bar b$ of $66.4\GeV$ (magenta star). All five emissions sit within $\Delta R = 0.31$ of the $b$: on this scale
the jet is already one blob.}
\label{fig:02_PS_display}
\end{figure}

\section{Open problems for Phase-2}

The shower of this chapter is steered by two python files of the release, both imported by the fragment
\path{Configuration/GenProduction/python/BBbar_14TeV_TuneCP5_cfi.py}. What follows is what they contain in
\code{CMSSW\_15\_0\_5}, and what it means for the HL-LHC.

\paragraph{CP5 showers: a 2017 tune of 13 TeV data, extrapolated.}
\path{Configuration/Generator/python/MCTunes2017/PythiaCP5Settings_cfi.py} starts from \code{Tune:pp 14} and
\code{Tune:ee 7} (the Monash 2013 tune of Pythia8); for the showers it changes only the coupling, setting
\code{SpaceShower:alphaSvalue} and \code{TimeShower:alphaSvalue} to 0.118 with \code{alphaSorder=2}. The
directory name says when: the CP tunes were derived in 2017 from underlying-event measurements at 13 TeV
(CMS-GEN-17-001). The energy dependence of the tune is the single parameter \code{MultipartonInteractions:ecmPow =
0.03344} in the formula $\pt^0 = \pt^{0,\mathrm{ref}}\,(E_\mathrm{cm}/E_\mathrm{cm}^\mathrm{ref})^{\mathrm{ecmPow}}$
with \code{pT0Ref = 1.41} and the reference energy at its Pythia8 default \code{MultipartonInteractions:ecmRef =
7000} GeV (not set in the file): the cut-off is $1.442\GeV$ at $13.6\TeV$ and $1.443\GeV$ at $14\TeV$, a change of
one per mille, so going from Run 3 to the HL-LHC energy is harmless for the formula; what is open is whether a
power law fitted between 7 and 13 TeV holds at 14, which no file in the release addresses. This parameter enters
chapter 04, not this one: for the ISR the regularisation is \code{SpaceShower:pT0Ref = 2.0} with
\code{SpaceShower:ecmPow = 0.0} and \code{samePTasMPI = off}, Pythia8 defaults that CMS never tuned, so the ISR
cut-off has no energy dependence at all. The FSR cut-off \code{TimeShower:pTmin = 0.5} is also untouched, and Pythia8
overrides it at initialisation (raised to $0.611\GeV$, in every log of this book). At the HL-LHC the soft end of
the ISR spectrum is what populates $1.5 < \abseta < 3.0$, the acceptance of the HGCal, with low-\pt{} partons on
top of 200 pileup collisions; the cut-off that decides how many of them exist is a default of 2014.

\paragraph{The ``Run 3'' copy of CP5 only changed two cross sections.} The release also has
\path{Configuration/Generator/python/MCTunesRun3ECM13p6TeV/PythiaCP5Settings_cfi.py}. Its \code{diff} against the
2017 file is two lines: \code{SigmaTotal:sigmaEl = 22.08} and \code{SigmaTotal:sigmaTot = 101.037} instead of
$21.89$ and $100.309$ (the total and elastic $pp$ cross sections used with \code{SigmaTotal:mode = 0}, which fixes
them by hand instead of letting Pythia8 compute them for the actual energy). Every shower, MPI and colour-reconnection
parameter is byte-identical. Our fragment, like all 40 \code{*14TeV*} fragments of
\path{Configuration/Generator/python/} that use CP5 (\path{TTbar_14TeV_TuneCP5_cfi.py},
\path{MinBias_14TeV_pythia8_TuneCP5_cfi.py}, \path{QCD_Pt_80_120_14TeV_TuneCP5_cfi.py}, \ldots), imports the
2017 file; none of them imports the Run 3 one (0 of 40), and there is no \code{MCTunes} directory for 14 TeV or
Phase-2. The Phase-2 pileup of the release validation is built from exactly these fragments:
\path{Configuration/PyReleaseValidation/python/relval_steps.py} uses \code{MinBias\_14TeV\_pythia8\_TuneCP5} with
\code{AVE\_200\_BX\_25ns} for the \code{Run4} datasets. So the 200 minimum-bias collisions of a Phase-2 event, and
the showers of its signal, run with the 13 TeV numbers of 2017 and even with the 13 TeV total cross section that the
Run 3 copy corrected.

\paragraph{The shower-uncertainty weights carry a note that says it.}
\path{Configuration/Generator/python/PSweightsPythia/PythiaPSweightsSettings_cfi.py}, the block that adds the
ISR/FSR $\mu_R$ variations as event weights (not included in our fragment, which is why the GEN file carries no
such weights), ends with \code{UncertaintyBands:ISRpTmin2Fac = 20} and the comment ``for consistency with UL and
P8.240 set to 20, to be optimized and changed for Run 3 re-MC''. The optimisation announced for Run 3 is not in
\code{CMSSW\_15\_0\_5}; for Phase-2 nothing is planned in the file.

\paragraph{Pythia8CommonSettings has no version at all.}
\path{Configuration/Generator/python/Pythia8CommonSettings_cfi.py} holds eight lines used by every CMS Pythia8
fragment of Run 2, Run 3 and Phase-2 alike: \code{Beams:setProductionScalesFromLHEF} = off,
\code{Tune:preferLHAPDF} = 2, \code{Main:timesAllowErrors} = 10000, \code{Check:epTolErr} = 0.01 (the energy--momentum
conservation tolerance, loosened by a factor 100 from the Pythia8 default $10^{-4}$, visible in the settings table
of the log), \code{SLHA:minMassSM} = 1000, \code{ParticleDecays:limitTau0} = on, \code{ParticleDecays:tau0Max} = 10 and
\code{ParticleDecays:allowPhotonRadiation} = on. None of them concerns the shower directly; \code{tau0Max = 10}
(in mm) will matter in chapters 06 and 09, because it decides which particles decay in Pythia8 and which are handed
to Geant4, and in Phase-2 the first tracker layer and the MTD timing layer sit where that boundary falls. The file
carries no date and \code{cvmfs} carries no git history, so I could not establish from the release when it was last
changed; its content is the same for the 13 TeV and the 14 TeV fragments in the same directory.

\paragraph{What the event itself says.} In this event the $\bar b$ at $\eta = -3.75$ radiated nothing, by chance.
On average it radiates, and its gluons land between $\abseta = 3$ and $5$: in the 2024 detector this is the HF, in
Phase-2 it is the HF too, just outside the HGCal. The $b$ at $\eta = -1.44$ radiated five partons within $\Delta R <
0.31$; their number and their energy sharing are what a jet algorithm, a $b$ tagger or a jet-mass measurement sees,
and they are governed by $\alpha_s^\mathrm{FSR} = 0.118$ and \code{TimeShower:pTmin} of
table~\ref{tab:02_PS_settings}. Neither has been revisited for 14 TeV, and chapters 04 and 07 will show how much of
the jet they make.
